\documentclass[11pt]{article}
\usepackage{docstyle}
\renewcommand\sheetname{Homework 2 \textperiodcentered{} Metallic materials}
\begin{document}
\sheethead{Homework 2: Metallic materials}{Year 9 Science \textperiodcentered{} about 40 minutes}{Section A: A \quad Section B: M / E}

\section{Section A}
\Q Write the meaning of each word.\lvl{A}
\par\sub{a} density
\alines{1}{how much mass is in a given volume (mass divided by volume)}
\sub{b} alloy
\alines{1}{a mixture of two or more elements, at least one of which is a metal}
\sub{c} corrosion
\alines{1}{a chemical reaction of a metal with substances around it, such as oxygen and water}

\Q Name the heat treatment.\lvl{A}
\par\sub{a} A blacksmith wants a horseshoe to be soft, so it is easy to shape. \hfill\blank[34mm]{annealing}
\par\sub{b} The edge of a chisel must be very hard. \hfill\blank[34mm]{quenching}\par\vspace{1mm}

\Q Write the word equations.\lvl{A}
\par\sub{a} lead + oxygen $\longrightarrow$ \blank[70mm]{lead oxide}
\par\sub{b} sodium + hydrochloric acid $\longrightarrow$ \blank[70mm]{sodium chloride + hydrogen}\par\vspace{1mm}

\section{Quick review}
\Q Name the property of metals that makes each one a good choice.\lvl{A}
\par\sub{a} the copper base of a saucepan \hfill\blank[50mm]{good conductor of heat}
\par\sub{b} gold contacts inside a phone \hfill\blank[50mm]{conducts electricity, unreactive}
\par\sub{c} a steel bridge cable \hfill\blank[50mm]{strong}\par\vspace{1mm}

\section{Section B}
\Q A gold-coloured crown has a mass of 965\,g and a volume of 50.0\,cm$^{3}$. The density of pure gold is 19.3\,g/cm$^{3}$. Is the crown pure gold? Show your working.\lvl{M}
\flines{2}{density $= \dfrac{965\ \mathrm{g}}{50.0\ \mathrm{cm^{3}}} = 19.3$ g/cm$^{3}$\quad This matches pure gold, so it could be pure gold.}

\Q Aircraft bodies are made of an aluminium alloy, not steel. Use the data to explain why.\lvl{E}
\par\vspace{1mm}
\begin{tabular}{@{}l c c@{}}
\toprule
Material & Density (g/cm$^{3}$) & Strength\\
\midrule
aluminium alloy & 2.8 & high\\
steel & 7.8 & very high\\
\bottomrule
\end{tabular}
\alines{4}{The aluminium alloy is strong enough, and its density (2.8) is much lower than steel's (7.8), almost three times less. A lighter aircraft needs less fuel to fly. Steel is stronger, but it would make the aircraft far too heavy.}

\Q The steel exhaust pipe of a car rusts faster than the steel body of the car. Explain why.\lvl{M}
\alines{3}{The exhaust pipe gets hot, and the burning petrol makes water, so the pipe has water, oxygen and a high temperature. A raised temperature speeds up rusting.}

\Q A student adds the same mass of four metals to the same volume of dilute acid and measures the temperature rise.\lvl{E}
\par\vspace{1mm}
\begin{tabular}{@{}l c c c c@{}}
\toprule
Metal & magnesium & zinc & iron & copper\\
Temperature rise (\textdegree C) & 18 & 7 & 3 & 0\\
\bottomrule
\end{tabular}\par\vspace{1mm}
\sub{a} Put the metals in order of reactivity, most reactive first.
\alines{1}{magnesium, zinc, iron, copper}
\sub{b} Use the data to explain your order, and name two things the student kept the same.
\alines{4}{The more reactive the metal, the faster it reacts with the acid, so it gives out more heat: magnesium raises the temperature most (18\,\textdegree C) and copper does not react at all (0\,\textdegree C). Kept the same: mass of metal, volume of acid (also concentration of acid, starting temperature).}

\Q Sam writes about alloys and reactions.\lvl{E}
\par\vspace{1mm}
\samcard{0.9\linewidth}{\flawed{1. Steel is harder than pure iron because all the atoms in steel are the same size, so the layers slide easily.\\
2. magnesium + hydrochloric acid makes magnesium chloride + water\\
3. Bronze is an alloy, so it is a mixture of metals.\\
4. When copper is added to dilute hydrochloric acid, bubbles of hydrogen are given off.}}\par\vspace{1mm}
\textbf{Sam's answer has 3 errors.} Find each one, fix it, and say why it loses marks.
\alines{5}{1. Steel has atoms of different sizes (carbon), which disrupt the layers so they cannot slide easily (\Lref{alloy}). 2. Metal + acid gives hydrogen, not water: magnesium chloride + hydrogen (\Lref{wordeq}). 3. Copper is below the acids in the reactivity series, so there are no bubbles: no reaction (\Lref{react}).}

\vspace{4mm}
{\small Also set: Metals booklet p.10--11 (draw the solder graph, then answer Q1--2) and the crossword on p.16.}
\end{document}
